% Same frozen perfect-trained posteriors; inference only.
% Numerical evidence: assets/c1_evaluation.json and tables/c1_evaluation.csv.
\paragraph{Two evaluations of the perfect-training guarantee.}
We freeze the same three perfect-trained two-moons posteriors, with $n=512$
and hidden width eight. Both deployment cases retain the trained weights,
Gaussian prior, noiseless predictor, empirical training risk, and data
penalty. Only the channel-shift evaluation changes.

For independent feature erasures with probability $p$, choose the discrete
metric $d_{\mathcal C}(c,c')=\mathbf1\{c\ne c'\}$ on complete masks.
Every $[0,1]$ loss is uniformly $1$-Lipschitz for this metric. Thus C1's
first inequality has the explicit expected-distance term
\begin{equation}
 B_{\rm erasure}=\mathbb E_Pd_{\mathcal C}(C,c_0)=1-(1-p)^8.
 \label{equ:sim_c1_erasure}
\end{equation}
This equals the earlier total-variation specialization numerically.

For the second case, the entire trained network transmits its two output
logits, and the receiver applies softmax after additive, centered lognormal
interference. The deployment law is fully specified by independent
$Z_i\sim\mathcal N(0,1)$ and
\begin{equation}
 N_i=\frac{\sigma}{\sqrt{v_\tau}}
       (e^{\tau Z_i-\tau^2/2}-1),\quad
 v_\tau=e^{\tau^2}-1,\quad \tau=0.1.
 \label{equ:sim_c1_lognormal_channel}
\end{equation}
Each coordinate has zero mean and variance $\sigma^2$, but its density is
skewed and non-Gaussian. The fixed shape is chosen analytically before the
deployment-risk evaluation. Moving the communication split to the output
preserves exactly the noiseless trained predictor. This is an additive-output
case, rather than a Lipschitz claim for the earlier hidden-feature channel.
For the smooth bounded loss, $L=\log(1/\varepsilon_\ell)$ and
\begin{equation}
 \nabla_{\rm logits}\ell=
 \frac{(1-\varepsilon_\ell)p_y}
 {\varepsilon_\ell+(1-\varepsilon_\ell)p_y}
 \frac{\boldsymbol p-\boldsymbol e_y}{L},\qquad
 \|\nabla_{\rm logits}\ell\|_2\leq\frac{\sqrt2}{L}.
 \label{equ:sim_c1_logit_lipschitz}
\end{equation}
Consequently, $K=\sqrt2/\log100$ is uniform over all weights and inputs for
the Euclidean noise metric. This argument uses the smooth loss, rather
than the discontinuous $0$--$1$ loss.

The first expected-distance term now requires integration of
\begin{equation}
 K\mathbb E_P\|N\|_2=
 \frac{K\sigma}{\sqrt{v_\tau}}\,
 \mathbb E\sqrt{\sum_{i=1}^2(e^{\tau Z_i-\tau^2/2}-1)^2}.
 \label{equ:sim_c1_lognormal_distance}
\end{equation}
This fully known channel has no simple known closed-form expression for
its mean Euclidean distance. A direct evaluation therefore requires
numerical integration or Monte Carlo. We retain an independent two-million-draw
MC estimate as a diagnostic, with its standard error in the source data;
it does not enter the analytic guarantee.

Instead choose the fixed, moment-matched Gaussian reference
$Q_\sigma=\mathcal N(0,\sigma^2 I_2)$. Its support contains that of $P$.
Writing $a=\sigma/\sqrt{v_\tau}$, the coordinate entropy is
$h(N_i)=\log(a\tau)-\tau^2/2+\tfrac12\log(2\pi e)$.
Together with $\mathbb E N_i^2=\sigma^2$, this yields the exact KL
\begin{equation}
 D(P\Vert Q_\sigma)
 =J_\tau=\tau^2+\log\frac{e^{\tau^2}-1}{\tau^2}.
 \label{equ:sim_c1_lognormal_kl}
\end{equation}
C1's second inequality is therefore an entirely closed-form evaluation:
\begin{align}
 K\mathbb E_P\|N\|_2
 &\leq B_{\rm ref}(\sigma)=\log M_2(K\sigma)+J_\tau,\nonumber\\
 M_2(t)&=1+t\sqrt{2\pi}\,e^{t^2/2}\Phi(t),
 \label{equ:sim_c1_reference}
\end{align}
where $\Phi$ is the standard normal CDF. Neither deployment-distance
integration nor numerical reference optimization is needed. For positive
$\sigma$ and $\tau$, the lognormal tail gives
$\mathbb E_P e^{K\|N\|_2}=\infty$, so choosing $Q=P$ would not be useful;
the Gaussian reference has the required finite exponential moment. The first
inequality remains valid; the second provides a closed-form relaxation
when the first expectation has no simple analytic expression.
A further control is the elementary Jensen upper bound
$K\mathbb E\|N\|_2\leq K\sigma\sqrt2$. Neither analytic upper bound
uniformly dominates the other. At $\sigma=0$, both channel laws are the
identity and their shift penalties are exactly zero; this endpoint is
treated separately from the continuous Gaussian reference, whose positive-noise
KL term retains the non-Gaussian shape penalty $J_\tau$.

For each posterior we add its empirical risk, a fresh 100-draw empirical-MC
upper correction, and the unweighted data penalty. The conservative data-event
budget is $0.04/473$, covering the 473 previously explored journal trials;
the three posteriors are fixed before the fresh MC evaluations, which each
receive budget $0.01/3$. Thus the declared family has failure probability
at most $0.05$. Analytic channel-shift bounds cover every displayed severity.
Test risks remain finite-test estimates.

Figure~\ref{fig:sim_c1_evaluation} compares the two evaluations;
Table~\ref{tab:c1_evaluation} in the supplement gives numerical values.
\begin{figure*}[t]
 \centering
 \includegraphics[width=0.92\textwidth]{\simroot/assets/fig_c1_evaluation.pdf}
 \caption{C1 under identical perfect-channel training. Feature erasure uses
 the exact first expected-distance term. Centered lognormal output noise
 uses the closed-form Gaussian-reference expression in the second inequality.
 Risk-upper curves include empirical risk, its MC correction, and the data
 penalty. The lognormal direct-expectation MC curve is a diagnostic, not a
 certified risk upper; its distance estimates never enter the reference bound.
 The Jensen curve is a separate analytic upper bound. Held-out-risk error bars
 show one sample SD over all three seeds. Uncapped bounds above one are vacuous.}
 \label{fig:sim_c1_evaluation}
\end{figure*}

The shared mean empirical risk, data penalty, and empirical-MC correction
are $0.0746$, $0.4240$, and $0.1689$. At $\sigma=0.5$,
$J_\tau=0.0150042$ and $\log M_2(K\sigma)=0.1976094$, giving the
closed-form second shift upper $0.2126135$ and corrected mean risk upper
$0.8800855$. The Jensen shift upper is $0.2171472$. This example exhibits
the intended tradeoff: the first expectation needs numerical evaluation,
whereas a chosen reference produces an explicit, looser analytic guarantee
without changing training.
